\section{Répétitions}
%\label{sec:7:1:repetitions}

Nous avons déjà vu quelques exemples de programmes \python\ dans lesquels une
séquence de commandes du \robot\ ou d'appels de procédures se trouve répétée de
nombreuses fois de suite. C'est le cas par exemple dans le programme {\idt FriseDeN\_1}
dont l'exécution produit la frise ci-dessous~:

%\vspace{\baselineskip}
\vspace*{8mm}%
\noindent\includevspacePostScript[1cm]{FriseDeN}{0.95}{0.95}
\vspace*{-1cm}%
{
\parindent \parindentprog
\medskip

\carrenoir
\vspace{\neglskip}\mbox{\blc {\mcl from}{\idt \ LibRobot\ }{\mcl import}{\idt \ *}}\par
\vspace{0.3\baselineskip}
\vspace{\neglskip}\mbox{\blc {\cmt \#\ Dessiner\ une\ frise\ de\ N}}\par
\vspace{\neglskip}\mbox{\blc {\idt ag()}}\par
\vspace{\neglskip}\mbox{\blc {\idt av();\ av();\ pdt();\ vg();\ vd();\ pdt()}}\par
\vspace{\neglskip}\mbox{\blc {\idt av();\ av();\ pdt();\ vg();\ vd();\ pdt()}}\par
\vspace{\neglskip}\mbox{\blc {\idt av();\ av();\ pdt();\ vg();\ vd();\ pdt()}}\par
\vspace{\neglskip}\mbox{\blc {\idt av();\ av();\ pdt();\ vg();\ vd();\ pdt()}}\par
\vspace{\neglskip}\mbox{\blc {\idt av();\ av();\ pdt();\ vg();\ vd();\ pdt()}}\par
\vspace{\neglskip}\mbox{\blc {\idt av();\ av();\ pdt();\ vg();\ vd();\ pdt()}}\par
\vspace{\neglskip}\mbox{\blc {\idt av();\ av();\ pdt();\ vg();\ vd();\ pdt()}}\par
\vspace{\neglskip}\mbox{\blc {\idt av();\ av();\ pdt();\ vg();\ vd();\ pdt()}}\par
\vspace{\neglskip}\mbox{\blc {\idt av();\ av();\ pdt();\ vg();\ vd();\ pdt()}}\par
\vspace{\neglskip}\mbox{\blc {\idt av();\ av();\ pdt();\ vg();\ vd();\ pdt()}}\par
\vspace{\neglskip}\mbox{\blc {\idt av();\ av();\ pdt();\ vg();\ vd();\ pdt()}}\par
\vspace{\neglskip}\mbox{\blc {\idt av();\ av();\ pdt();\ vg();\ vd();\ pdt()}}\par
\vspace{\neglskip}\mbox{\blc {\idt av();\ av();\ pdt();\ vg();\ vd();\ pdt()}}\par
\vspace{\neglskip}\mbox{\blc {\idt av();\ av();\ pdt();\ vg();\ vd();\ pdt()}}\par
\vspace{\neglskip}\mbox{\blc {\idt av();\ av();\ pdt();\ vg();\ vd();\ pdt()}}\par
\vspace{\neglskip}\mbox{\blc {\idt av();\ av();\ pdt();\ vg();\ vd();\ pdt()}}\par
\vspace{\neglskip}\mbox{\blc {\idt st()}}\par

\medskip }


Pour éviter les répétitions successives des mêmes séquences d'ins\-truc\-tions,
nous allons étudier une nouvelle instruction du langage \python\ qui s'appelle la
{\em boucle ``for''}\index{Boucle ``for''} (le mot anglais ``for'', qui signifie ``pour'',
se prononce comme il s'écrit%
\footnote{\ Nous parlons de boucle ``for'' et non pas de boucle ``pour'' pour rester dans le
contexte du vocabulaire utilisé dans le programme.}). Voici le
programme {\idt FriseDeN\_1} de dessin de la frise ci-dessous
réécrit en utilisant une boucle ``for''~:

%\vspace{\baselineskip}
\vspace*{8mm}%
\noindent\includevspacePostScript[1cm]{FriseDeNprime}{0.95}{0.95}
\vspace*{-1cm}%
\label{prog:FriseDeN_2}
{
\parindent \parindentprog
\medskip

\carrenoir
\vspace{\neglskip}\mbox{\blc {\mcl from}{\idt \ LibRobot\ }{\mcl import}{\idt \ *}}\par
\vspace{0.3\baselineskip}
\vspace{\neglskip}\mbox{\blc {\cmt \#\ Dessiner\ une\ frise\ de\ N\ avec\ une\ boucle\ \char34{}for\char34{}}}\par
\vspace{\neglskip}\mbox{\blc {\idt ag()}}\par
\vspace{\neglskip}\mbox{\blc {\mcl for}{\idt \ j\ }{\mcl in}{\idt \ range(16):}}\par
\vspace{\neglskip}\mbox{\blc \blk\blk{\idt av();\ av();\ pdt();\ vg();\ vd();\ pdt()}}\par
\vspace{\neglskip}\mbox{\blc {\idt st()}}\par

\medskip }


\noindent En français, la boucle `{\mcl for}\ {\idt j}\ {\mcl in}\ {\idt range(16):}\ \ldots'
peut se traduire par  `pour j parmi les 16 premiers nombres faire \ \ldots'. Le mot anglais {\mcl in},
qui signifie `dans', se prononce 'inne'. Le mot anglais {\idt range}, qui signifie
`suite' ou `étendue', se prononce 'rèïndje'.

La séquence d'instructions à répéter s'écrit sur les lignes qui suivent les deux-points `{\idt :}',
décalées vers la droite. Elle s'appelle le
{\em corps de la boucle}\index{Corps de boucle}. Le \index{Itération}{\em nombre d'itérations} c'est-à-dire le
nombre de fois qu'il faut répéter le corps de la boucle est indiqué
entre les parenthèses de {\idt range}. Dans notre exemple, il faut répéter 16 fois
de suite le corps de la boucle  `{\idt  av();\ av();\ pdt();\ vg();\ vd();\ pdt()}'.

\section{Le compteur de boucle}
%\label{sec:7:3:decls}

Juste après le mot {\mcl for} nous trouvons le nom {\idt j} d'un
{\em compteur}\index{Comp\-teur de bou\-cle}. Ce compteur va servir à l'ordinateur pour
compter le nombre de fois qu'il aura à répéter la séquence d'instructions dans la
boucle ``for'', exactement comme nous pourrions le faire en comptant sur nos  doigts$\:$!
Pour cela l'ordinateur utilise une case de sa mémoire, qui contient successivement
les nombres entiers 0, 1, 2, \ldots, 15. En effet, {\idt range(16)} désigne la suite des
16 nombres entiers qui commence à 0 et s'arrête juste avant 16 (comme les informaticiens,
\python\ compte à partir de 0). On dit que {\idt j} est une
{\em variable}\index{Variable}\index{Variable entière} entière puisque sa valeur varie
d'une itération à l'autre.

Nous devons donner un nom à ce compteur. C'est un identificateur que nous pouvons choisir
comme il nous plaît (ici par exemple {\idt  j}), en respectant la règle des
identificateurs~: il commence par une lettre (ou un souligné `\_'), ne contient que des lettres, des chiffres et des soulignés, et ne doit pas être un mot réservé de
\python. Si l'on écrit `{\mcl for}\ {\idt 1er}\ {\mcl in}\ \ldots', \python\ répond
`{\tt SyntaxError: invalid decimal literal}' (nombre décimal incorrect, puisque
{\idt 1er} commence par un chiffre). Si l'on écrit `{\mcl for}\ {\mcl in}\ {\mcl in}\ \ldots',
il répond `{\tt SyntaxError: invalid syntax}' (syntaxe incorrecte).

En \python, il n'est pas nécessaire d'annoncer à l'avance que l'on va utiliser
un compteur%
\footnote{\ D'autres langages comme \pascal\ obligent à {\em déclarer} le compteur au début
du programme ou de la procédure, en indiquant qu'il désigne un nombre entier.}~:
le compteur est créé par la boucle ``for'' elle-même. \python\
sait que c'est un nombre entier puisqu'il prend les valeurs 0, 1, 2, \ldots\ Après la
boucle, le compteur garde la dernière valeur qu'il a prise (15 dans notre exemple), mais
il vaut mieux ne pas s'en servir~: si la boucle n'a fait aucune itération, le
compteur n'existe pas et \python\ indique l'erreur
`{\tt NameError: name 'j' is not defined}' (le nom {\idt j} n'est pas défini).

Pour que le compteur prenne les valeurs 1, 2, 3, \ldots, 16, on écrirait
`{\mcl for}\ {\idt j}\ {\mcl in}\ {\idt range(1,\ 17):}'. De manière générale,
{\idt range(}$a${\idt ,\ }$b${\idt )} désigne la suite des nombres entiers de $a$ à $b-1$~:
elle commence à $a$ et s'arrête juste {\em avant} $b$. Ainsi {\idt range(1,\ 5)} désigne
1, 2, 3, 4 et {\idt range(}$n${\idt )}, qui est une abréviation de
{\idt range(0,\ }$n${\idt )}, désigne les $n$ nombres 0, 1, \ldots, $n-1$.
Pour que le compteur aille de 1 à $n$, on écrit donc
`{\mcl for}\ {\idt i}\ {\mcl in}\ {\idt range(1,\ }$n${\idt \ +\ 1):}'.

Le compteur de la boucle du programme {\idt FriseDeN\_2} ne sert qu'à compter les itérations.
Il est alors plus simple d'écrire {\idt range(16)}.
Nous verrons au chapitre suivant que l'on peut aussi utiliser sa valeur dans le corps de
la boucle.

\section{La boucle ``for''}
%\label{sec:7:4:for}

Une {\em boucle ``for''}\index{Boucle ``for''} a la forme suivante~:\\
\begin{center}
\begin{tabular}{@{}|p{12.5cm}|@{}} \hline
{\em Boucle ``for''~:}\\ \hline \hline
\begin{minipage}[t]{12.5cm}
{\medskip
\mbox{\blc\ \ \ {\mcl for}\hspace{0.3em}\blanc\hspace{0.3em}{\small\rm Identificateur}\hspace{0.3em}\blanc\hspace{0.3em}{\mcl in}\hspace{0.3em}\blanc\hspace{0.3em}{\idt range(}{\small\rm Nombre d'itérations}{\idt ):}}\par
\mbox{\blc\ \ \ \ \ \ \ \small\rm  \ldots\ Instructions du corps de la boucle (décalées) \ldots}\par
\medskip}
\end{minipage}\\ \hline
\end{tabular}
\end{center}
\noindent L'identificateur est le nom du compteur d'itérations.
La  séquence d'instructions qui est répétée s'appelle le {\em corps de
la boucle}\index{Corps de boucle}. Comme le corps d'une procédure, elle s'écrit sur les lignes
qui suivent les deux-points `{\idt :}', toutes décalées vers la droite du même nombre de
blancs (quatre en général), une instruction par ligne ou plusieurs instructions sur une même
ligne séparées par des `{\idt ;}'. C'est ce décalage, appelé {\em indentation}\index{Indentation},
qui indique à \python\ où commence et où finit le corps de la boucle~: la boucle se
termine à la première ligne qui n'est plus décalée.

Il ne faut pas oublier les blancs \hspace{0.3em}\blanc\hspace{0.3em} autour du mot {\mcl in}.
Si l'on écrit {\idt jin} au lieu de {\idt j}\ {\mcl in}, ou {\idt inrange} au lieu de
{\mcl in}\ {\idt range}, \python\ ne reconnaît plus le mot {\mcl in} et signale
`{\tt SyntaxError: invalid syntax}'. Si l'on oublie les deux-points `{\idt :}' à la fin
de la ligne du {\mcl for}, \python\ indique `{\tt SyntaxError: expected ':'}'
(deux-points attendu).

Une boucle ``for'' est elle-même une instruction. Les instructions qui la suivent
s'écrivent sans décalage, au même niveau que le {\mcl for}, comme le {\idt st()} du
programme {\idt FriseDeN\_2}. Si l'on décalait aussi la ligne {\idt st()}, elle ferait partie
du corps de la boucle et serait répétée 16 fois~! Si l'on oublie de décaler le corps
de la boucle, \python\ indique
`{\tt IndentationError: expected an indented block after 'for' statement on line 5}'
(erreur d'indentation~: un bloc décalé est attendu après l'instruction {\mcl for} de la
ligne 5). Si les lignes du corps ne sont pas toutes décalées du même nombre de blancs,
\python\ signale `{\tt IndentationError: unexpected indent}' (décalage inattendu) ou
`{\tt IndentationError: unindent does not match any outer indentation level}'
(le retour à gauche ne correspond à aucun niveau de décalage précédent).

Si le `Nombre d'itérations' est plus petit que 1 (0 par exemple)  les instructions du
corps de la boucle ne sont pas exécutées du tout. \underline{Il ne faut pas mettre} \underline{de point-virgule `{\idt $;$}'} \underline{après les deux-points}~: \python\ signale alors `{\tt SyntaxError: invalid syntax}' juste après les deux-points.

Quand le corps de la boucle est court, on peut écrire toute la boucle sur une seule
ligne, comme dans `{\mcl for}\ {\idt i\ }{\mcl in}\ {\idt range(10):\ av()}'.
Les blancs supplémentaires à l'intérieur d'une ligne servent seulement à rendre le
programme plus lisible, mais ceux du début de ligne sont importants.

\begin{exemple}[Frise de pointes]{\rm
 La frise suivante~:\\[20mm]
\includevspacePostScript[1cm]{FriseDePointes}{0.95}{0.95}%
\\[1ex]
est produite par le programme ci-dessous~:
{
\parindent \parindentprog
\medskip

\carrenoir
\vspace{\neglskip}\mbox{\blc {\mcl from}{\idt \ LibRobot\ }{\mcl import}{\idt \ *}}\par
\vspace{0.3\baselineskip}
\vspace{\neglskip}\mbox{\blc {\cmt \#\ Dessiner\ une\ frise\ de\ pointes}}\par
\vspace{\neglskip}\mbox{\blc {\idt ag();\ pqtd()}}\par
\vspace{\neglskip}\mbox{\blc {\mcl for}{\idt \ j\ }{\mcl in}{\idt \ range(32):}}\par
\vspace{\neglskip}\mbox{\blc \blk\blk{\idt lc();\ av();\ bc();\ pdt();\ vd();\ pdt()}}\par
\vspace{\neglskip}\mbox{\blc \blk\blk{\idt av();\ av();\ pdt();\ vd()}}\par
\vspace{\neglskip}\mbox{\blc {\idt st()}}\par

\medskip }

\noindent $\Box$}
\end{exemple}

\section{Utilisation de la boucle ``for'' dans une procédure}
%\label{sec:7:5:for-proc}

Le corps d'une procédure peut contenir des boucles ``for'', comme n'importe quelle
autre instruction~:
\index{Compteur de boucle}
\begin{center}
\begin{tabular}{@{}|p{12.5cm}|@{}} \hline
{\em Déclaration de procédure~:}\\ \hline \hline
\begin{minipage}[t]{12.5cm}
\medskip
\mbox{\blc\ \ \ {\mcl def}\hspace{0.3em}\blanc\hspace{0.3em}\idt nomDeLaProcedure():}\par
\mbox{\blc\ \ \ \ \ \ \ \small\rm \ldots\ Instructions du corps de la procédure%
\footnote{\ Y compris des boucles ``for'' utilisant des compteurs {\idt i}, {\idt j}, \ldots}
(décalées) \ldots}\par
\medskip
\end{minipage}\\ \hline
\end{tabular}
\end{center}

Puisque le compteur d'une boucle ne sert que dans le corps de cette boucle et qu'il
repart de sa première valeur à chaque nouvelle boucle, un même
identificateur peut servir de compteur dans plusieurs boucles différentes,
qu'elles soient dans des procédures différentes ou les unes après les autres
dans la même procédure. Ceci est illustré dans le programme
{\idt Fleche} ci-dessous où des compteurs de boucles
différents, utilisés dans des procédures différentes ({\idt
empennage} et {\idt tige}), portent le même nom {\idt i}.

\begin{exemple}[Flèche]{\rm \ \\[7cm]
\includePostScript{FlecheVerticale}{0.95}{0.95}%
\vspace*{-7cm}%
{
\parindent \parindentprog
\medskip

\carrenoir
\vspace{\neglskip}\mbox{\blc {\mcl from}{\idt \ LibRobot\ }{\mcl import}{\idt \ *}}\par
\vspace{0.3\baselineskip}
\vspace{\neglskip}\mbox{\blc {\cmt \#\ Dessiner\ une\ flèche\ verticale}}\par
\vspace{\neglskip}\mbox{\blc {\mcl def}{\idt \ plume():}}\par
\vspace{\neglskip}\mbox{\blc \blk\blk{\idt av();\ p3htg();\ av();\ lc();\ p3htg();\ av();\ av()}}\par
\vspace{\neglskip}\mbox{\blc \blk\blk{\idt p3htg();\ bc();\ av();\ phtd()}}\par
\vspace{0.3\baselineskip}
\vspace{\neglskip}\mbox{\blc {\mcl def}{\idt \ empennage():}}\par
\vspace{\neglskip}\mbox{\blc \blk\blk{\mcl for}{\idt \ i\ }{\mcl in}{\idt \ range(5):}}\par
\vspace{\neglskip}\mbox{\blc \blk\blk\blk\blk{\idt plume()}}\par
\vspace{0.3\baselineskip}
\vspace{\neglskip}\mbox{\blc {\mcl def}{\idt \ tige():}}\par
\vspace{\neglskip}\mbox{\blc \blk\blk{\mcl for}{\idt \ i\ }{\mcl in}{\idt \ range(10):}}\par
\vspace{\neglskip}\mbox{\blc \blk\blk\blk\blk{\idt av()}}\par
\vspace{0.3\baselineskip}
\vspace{\neglskip}\mbox{\blc {\mcl def}{\idt \ pointe():}}\par
\vspace{\neglskip}\mbox{\blc \blk\blk{\idt pqtg();\ av();\ p3htd();\ av();\ pqtd();\ av()}}\par
\vspace{\neglskip}\mbox{\blc \blk\blk{\idt p3htd();\ av()}}\par
\vspace{0.3\baselineskip}
\vspace{\neglskip}\mbox{\blc {\idt eb();\ empennage();\ tige();\ pointe();\ st()}}\par

\medskip }

\noindent $\Box$ }\end{exemple}


\section{Boucles ``for'' imbriquées (in\-dé\-pen\-dan\-tes)}
%\label{sec:7:6:for-for-ind}

Le corps d'une boucle ``for''~:

\mbox{\blc\ \ \ {\mcl for}\hspace{0.3em}\blanc\hspace{0.3em}{\small\rm Identificateur1}\hspace{0.3em}\blanc\hspace{0.3em}{\mcl in}\hspace{0.3em}\blanc\hspace{0.3em}{\idt range(}{\small\rm Nombre d'itérations1}{\idt ):}}\par
\mbox{\blc\ \ \ \ \ \ \ \small\rm  \ldots}\par

\noindent peut à nouveau contenir des boucles ``for''~:

\mbox{\blc\ \ \ {\mcl for}\hspace{0.3em}\blanc\hspace{0.3em}{\small\rm Identificateur1}\hspace{0.3em}\blanc\hspace{0.3em}{\mcl in}\hspace{0.3em}\blanc\hspace{0.3em}{\idt range(}{\small\rm Nombre d'itérations1}{\idt ):}}\par
\mbox{\blc\ \ \ \ \ \ \ \small\rm  \ldots}\par
\mbox{\blc\ \ \ \ \ \ \ {\mcl for}\hspace{0.3em}\blanc\hspace{0.3em}{\small\rm Identificateur2}\hspace{0.3em}\blanc\hspace{0.3em}{\mcl in}\hspace{0.3em}\blanc\hspace{0.3em}{\idt range(}{\small\rm Nombre d'itérations2}{\idt ):}}\par
\mbox{\blc\ \ \ \ \ \ \ \ \ \ \ \small\rm  \ldots}\par
\mbox{\blc\ \ \ \ \ \ \ \small\rm  \ldots}\par

\noindent On dit alors que les  boucles ``for'' sont {\em imbriquées}.
La boucle utilisant le
compteur {\small\rm Identificateur2} pour faire
{\small\rm Nombre d'itérations2} itérations s'appelle la
 {\em boucle intérieure}. La boucle utilisant le compteur {\small\rm Identificateur1} pour faire
{\small\rm Nombre d'itérations1} itérations s'appelle la {\em boucle extérieure}.
Le corps de la boucle intérieure est décalé deux fois~: une fois parce qu'il fait partie
du corps de la boucle extérieure, une deuxième fois parce qu'il est le corps de la boucle
intérieure. La boucle intérieure se termine à la première ligne qui revient au décalage
de son {\mcl for}~; la boucle extérieure à la première ligne qui revient au décalage du
premier {\mcl for}. C'est donc le décalage qui indique, sans ambiguïté, à quel {\mcl for}
appartient chaque instruction. Considérons par exemple
le dessin d'un carré~:
\par\noindent\begin{minipage}{\linewidth}
\vbox to 0pt{\includevspacePostScript[20mm]{Carre1}{0.95}{0.95}}%
{
\parindent \parindentprog
\medskip

\carrenoir
\vspace{\neglskip}\mbox{\blc {\mcl from}{\idt \ LibRobot\ }{\mcl import}{\idt \ *}}\par
\vspace{0.3\baselineskip}
\vspace{\neglskip}\mbox{\blc {\cmt \#\ Dessiner\ un\ carré\ avec\ des\ procédures}}\par
\vspace{0.3\baselineskip}
\vspace{\neglskip}\mbox{\blc {\cmt \#\ Dessiner\ un\ côté\ du\ carré\ et\ pivoter\ à\ droite}}\par
\vspace{\neglskip}\mbox{\blc {\mcl def}{\idt \ cote():}}\par
\vspace{\neglskip}\mbox{\blc \blk\blk{\mcl for}{\idt \ j\ }{\mcl in}{\idt \ range(4):}}\par
\vspace{\neglskip}\mbox{\blc \blk\blk\blk\blk{\idt av()}}\par
\vspace{\neglskip}\mbox{\blc \blk\blk{\idt pqtd()}}\par
\vspace{0.3\baselineskip}
\vspace{\neglskip}\mbox{\blc {\cmt \#\ Dessiner\ un\ carré}}\par
\vspace{\neglskip}\mbox{\blc {\mcl def}{\idt \ carre():}}\par
\vspace{\neglskip}\mbox{\blc \blk\blk{\mcl for}{\idt \ i\ }{\mcl in}{\idt \ range(4):}}\par
\vspace{\neglskip}\mbox{\blc \blk\blk\blk\blk{\idt cote()}}\par
\vspace{0.3\baselineskip}
\vspace{\neglskip}\mbox{\blc {\idt carre();\ st()}}\par

\medskip }

\end{minipage}

En remplaçant les appels des procédures par leurs corps, on obtient
deux boucles imbriquées~:

%\includePostScript{Carre1}{0.95}{0.95}
{
\parindent \parindentprog
\medskip

\carrenoir
\vspace{\neglskip}\mbox{\blc {\mcl from}{\idt \ LibRobot\ }{\mcl import}{\idt \ *}}\par
\vspace{0.3\baselineskip}
\vspace{\neglskip}\mbox{\blc {\cmt \#\ Dessiner\ un\ carré\ avec\ des\ boucles\ \char34{}for\char34{}\ imbriquées}}\par
\vspace{\neglskip}\mbox{\blc {\cmt \#\ Dessiner\ un\ carré}}\par
\vspace{\neglskip}\mbox{\blc {\mcl for}{\idt \ i\ }{\mcl in}{\idt \ range(4):}}\par
\vspace{\neglskip}\mbox{\blc \blk\blk{\cmt \#\ Dessiner\ un\ côté\ du\ carré\ et\ pivoter\ à\ droite}}\par
\vspace{\neglskip}\mbox{\blc \blk\blk{\mcl for}{\idt \ j\ }{\mcl in}{\idt \ range(4):}}\par
\vspace{\neglskip}\mbox{\blc \blk\blk\blk\blk{\idt av()}}\par
\vspace{\neglskip}\mbox{\blc \blk\blk{\idt pqtd()}}\par
\vspace{\neglskip}\mbox{\blc {\idt st()}}\par

\medskip }


La boucle intérieure sert à construire un côté tandis que la boucle
ex\-té\-rieu\-re dessine les quatre côtés.

\newpage
\section{Boucles ``for'' imbriquées (dépendantes)}
%\label{sec:7:7:for-for-dep}

Dans des boucles ``for'' imbriquées, le nombre d'itérations
dans la boucle in\-té\-rieu\-re peut être dif\-fé\-rent à chaque fois que le corps de
la boucle extérieure est exécutée. Commençons par considérer l'exemple
simple de segments de droite de longueurs 1, 2, 3, \ldots disposés en spirale
(le dessin est représenté ci-dessous à l'échelle $\frac{1}{2}$ avec la grille)~:\\[5cm]
\includePostScript{SpiraleCarree-1}{0.95}{0.95}%
\vspace*{-5cm}%
{
\parindent \parindentprog
\medskip

\carrenoir
\vspace{\neglskip}\mbox{\blc {\mcl from}{\idt \ LibRobot\ }{\mcl import}{\idt \ *}}\par
\vspace{0.3\baselineskip}
\vspace{\neglskip}\mbox{\blc {\cmt \#\ Dessiner\ une\ spirale\ carrée}}\par
\vspace{\neglskip}\mbox{\blc {\mcl for}{\idt \ i\ }{\mcl in}{\idt \ range(1,\ 23):}}\par
\vspace{\neglskip}\mbox{\blc \blk\blk{\mcl for}{\idt \ j\ }{\mcl in}{\idt \ range(i):}}\par
\vspace{\neglskip}\mbox{\blc \blk\blk\blk\blk{\idt av()}}\par
\vspace{\neglskip}\mbox{\blc \blk\blk{\idt pqtd()}}\par
\vspace{\neglskip}\mbox{\blc {\idt st()}}\par

\medskip }

\vspace{2.6cm}% le listing Python est plus court que la figure placee a cote

Le nombre d'itérations dans la boucle ``for'' intérieure est égal à la valeur
du compteur {\idt i} de la boucle extérieure, c'est-à-dire successivement à 1, 2,
3, \ldots, 22 (puisque {\idt range(1,\ 23)} s'arrête juste avant 23). Le programme {\idt SpiraleCarree\_1} est donc équivalent au
programme {\idt SpiraleCarree\_2} ci-dessous où les boucles ``for'' ont été
éliminées pour montrer exactement quelle est la séquence de commandes exécutée
par le \robot.

%\special{SpiraleCarree-1.ps 0.24 0.24 scale 0 -342 translate}
{
\parindent \parindentprog
\medskip

\carrenoir
\vspace{\neglskip}\mbox{\blc {\mcl from}{\idt \ LibRobot\ }{\mcl import}{\idt \ *}}\par
\vspace{0.3\baselineskip}
\vspace{\neglskip}\mbox{\blc {\cmt \#\ Dessiner\ une\ spirale\ carrée\ sans\ boucles\ \char34{}for\char34{}}}\par
\vspace{\neglskip}\mbox{\blc {\idt av();\ pqtd()}}\par
\vspace{\neglskip}\mbox{\blc {\idt av();\ av();\ pqtd()}}\par
\vspace{\neglskip}\mbox{\blc {\idt av();\ av();\ av();\ pqtd()}}\par
\vspace{\neglskip}\mbox{\blc {\idt av();\ av();\ av();\ av();\ pqtd()}}\par
\vspace{\neglskip}\mbox{\blc {\idt av();\ av();\ av();\ av();\ av();\ pqtd()}}\par
\vspace{\neglskip}\mbox{\blc {\idt av();\ av();\ av();\ av();\ av();\ av();\ pqtd()}}\par
\vspace{\neglskip}\mbox{\blc {\idt av();\ av();\ av();\ av();\ av();\ av();\ av()}}\par
\vspace{\neglskip}\mbox{\blc {\idt pqtd()}}\par
\vspace{\neglskip}\mbox{\blc {\idt av();\ av();\ av();\ av();\ av();\ av();\ av();\ av()}}\par
\vspace{\neglskip}\mbox{\blc {\idt pqtd()}}\par
\vspace{\neglskip}\mbox{\blc {\idt av();\ av();\ av();\ av();\ av();\ av();\ av();\ av()}}\par
\vspace{\neglskip}\mbox{\blc {\idt av();\ pqtd()}}\par
\vspace{\neglskip}\mbox{\blc {\idt av();\ av();\ av();\ av();\ av();\ av();\ av();\ av()}}\par
\vspace{\neglskip}\mbox{\blc {\idt av();\ av();\ pqtd()}}\par
\vspace{\neglskip}\mbox{\blc {\idt av();\ av();\ av();\ av();\ av();\ av();\ av();\ av()}}\par
\vspace{\neglskip}\mbox{\blc {\idt av();\ av();\ av();\ pqtd()}}\par
\vspace{\neglskip}\mbox{\blc {\idt av();\ av();\ av();\ av();\ av();\ av();\ av();\ av()}}\par
\vspace{\neglskip}\mbox{\blc {\idt av();\ av();\ av();\ av();\ pqtd()}}\par
\vspace{\neglskip}\mbox{\blc {\idt av();\ av();\ av();\ av();\ av();\ av();\ av();\ av()}}\par
\vspace{\neglskip}\mbox{\blc {\idt av();\ av();\ av();\ av();\ av();\ pqtd()}}\par
\vspace{\neglskip}\mbox{\blc {\idt av();\ av();\ av();\ av();\ av();\ av();\ av();\ av()}}\par
\vspace{\neglskip}\mbox{\blc {\idt av();\ av();\ av();\ av();\ av();\ av();\ pqtd()}}\par
\vspace{\neglskip}\mbox{\blc {\idt av();\ av();\ av();\ av();\ av();\ av();\ av();\ av()}}\par
\vspace{\neglskip}\mbox{\blc {\idt av();\ av();\ av();\ av();\ av();\ av();\ av()}}\par
\vspace{\neglskip}\mbox{\blc {\idt pqtd()}}\par
\vspace{\neglskip}\mbox{\blc {\idt av();\ av();\ av();\ av();\ av();\ av();\ av();\ av()}}\par
\vspace{\neglskip}\mbox{\blc {\idt av();\ av();\ av();\ av();\ av();\ av();\ av();\ av()}}\par
\vspace{\neglskip}\mbox{\blc {\idt pqtd()}}\par
\vspace{\neglskip}\mbox{\blc {\idt av();\ av();\ av();\ av();\ av();\ av();\ av();\ av()}}\par
\vspace{\neglskip}\mbox{\blc {\idt av();\ av();\ av();\ av();\ av();\ av();\ av();\ av()}}\par
\vspace{\neglskip}\mbox{\blc {\idt av();\ pqtd()}}\par
\vspace{\neglskip}\mbox{\blc {\idt av();\ av();\ av();\ av();\ av();\ av();\ av();\ av()}}\par
\vspace{\neglskip}\mbox{\blc {\idt av();\ av();\ av();\ av();\ av();\ av();\ av();\ av()}}\par
\vspace{\neglskip}\mbox{\blc {\idt av();\ av();\ pqtd()}}\par
\vspace{\neglskip}\mbox{\blc {\idt av();\ av();\ av();\ av();\ av();\ av();\ av();\ av()}}\par
\vspace{\neglskip}\mbox{\blc {\idt av();\ av();\ av();\ av();\ av();\ av();\ av();\ av()}}\par
\vspace{\neglskip}\mbox{\blc {\idt av();\ av();\ av();\ pqtd()}}\par
\vspace{\neglskip}\mbox{\blc {\idt av();\ av();\ av();\ av();\ av();\ av();\ av();\ av()}}\par
\vspace{\neglskip}\mbox{\blc {\idt av();\ av();\ av();\ av();\ av();\ av();\ av();\ av()}}\par
\vspace{\neglskip}\mbox{\blc {\idt av();\ av();\ av();\ av();\ pqtd()}}\par
\vspace{\neglskip}\mbox{\blc {\idt av();\ av();\ av();\ av();\ av();\ av();\ av();\ av()}}\par
\vspace{\neglskip}\mbox{\blc {\idt av();\ av();\ av();\ av();\ av();\ av();\ av();\ av()}}\par
\vspace{\neglskip}\mbox{\blc {\idt av();\ av();\ av();\ av();\ av();\ pqtd()}}\par
\vspace{\neglskip}\mbox{\blc {\idt av();\ av();\ av();\ av();\ av();\ av();\ av();\ av()}}\par
\vspace{\neglskip}\mbox{\blc {\idt av();\ av();\ av();\ av();\ av();\ av();\ av();\ av()}}\par
\vspace{\neglskip}\mbox{\blc {\idt av();\ av();\ av();\ av();\ av();\ av();\ pqtd()}}\par
\vspace{\neglskip}\mbox{\blc {\idt st()}}\par

\medskip }


\section{Invariants de boucles ``for''}
%\label{sec:7:8:inv}

Quand on dessine une frise comme celle-ci~:\\[9mm]
\includevspacePostScript[9mm]{FriseDeTildes}{0.95}{0.95}%
\vspace*{-9mm}%

\noindent avec une boucle ``for'', il faut que le \robot\ se trouve toujours dans
la même direction
et dans la même position par rapport au prochain motif qui va être
dessiné en exécutant
le corps de la boucle. Autrement dit, à la fin du corps de la
boucle il faut
remettre le \robot\ dans la position de départ~:
{
\parindent \parindentprog
\medskip

\carrenoir
\vspace{\neglskip}\mbox{\blc {\mcl from}{\idt \ LibRobot\ }{\mcl import}{\idt \ *}}\par
\vspace{0.3\baselineskip}
\vspace{\neglskip}\mbox{\blc {\cmt \#\ Dessiner\ une\ frise\ de\ tildes}}\par
\vspace{0.3\baselineskip}
\vspace{\neglskip}\mbox{\blc {\cmt \#\ Dessiner\ un\ motif\ en\ forme\ de\ tilde\ \char126{}}}\par
\vspace{\neglskip}\mbox{\blc {\mcl def}{\idt \ tilde():}}\par
\vspace{\neglskip}\mbox{\blc \blk\blk{\idt phtd();\ av();\ phtd();\ av()}}\par
\vspace{\neglskip}\mbox{\blc \blk\blk{\idt pqtd();\ av();\ av();\ pqtg()}}\par
\vspace{\neglskip}\mbox{\blc \blk\blk{\idt av();\ phtg();\ av()}}\par
\vspace{0.3\baselineskip}
\vspace{\neglskip}\mbox{\blc {\idt ag()}}\par
\vspace{\neglskip}\mbox{\blc {\mcl for}{\idt \ i\ }{\mcl in}{\idt \ range(10):}}\par
\vspace{\neglskip}\mbox{\blc \blk\blk{\idt tilde()}}\par
\vspace{\neglskip}\mbox{\blc \blk\blk{\cmt \#\ Reculer\ au\ début\ du\ tilde\ suivant}}\par
\vspace{\neglskip}\mbox{\blc \blk\blk{\idt lc();\ p3htg();\ av();\ bc()}}\par
\vspace{\neglskip}\mbox{\blc \blk\blk{\cmt \#\ Orienter\ le\ robot\ au\ nord}}\par
\vspace{\neglskip}\mbox{\blc \blk\blk{\idt pqtd()}}\par
\vspace{\neglskip}\mbox{\blc {\idt st()}}\par

\medskip }

On dit que la position du \robot\ est un {\em invariant}\index{Invariant} de la boucle en ce sens
qu'elle doit être la même à chaque itération.

Si dans notre exemple on dessine un motif puis le suivant sans remettre entre-temps
le \robot\ dans la position de départ, on obtient le résultat suivant, dont
il faut avouer qu'il n'est pas très ressemblant à l'original~:\\[5cm]
\includePostScript{FriseDeTildesErr1}{0.95}{0.95}%
\vspace*{-5cm}
{
\parindent \parindentprog
\medskip

\carrenoir
\vspace{\neglskip}\mbox{\blc {\mcl from}{\idt \ LibRobot\ }{\mcl import}{\idt \ *}}\par
\vspace{0.3\baselineskip}
\vspace{\neglskip}\mbox{\blc {\cmt \#\ Dessiner\ une\ frise\ de\ tildes\ (première\ erreur)}}\par
\vspace{0.3\baselineskip}
\vspace{\neglskip}\mbox{\blc {\cmt \#\ Dessiner\ un\ motif\ en\ forme\ de\ tilde\ \char126{}}}\par
\vspace{\neglskip}\mbox{\blc {\mcl def}{\idt \ tilde():}}\par
\vspace{\neglskip}\mbox{\blc \blk\blk{\idt phtd();\ av();\ phtd();\ av()}}\par
\vspace{\neglskip}\mbox{\blc \blk\blk{\idt pqtd();\ av();\ av();\ pqtg()}}\par
\vspace{\neglskip}\mbox{\blc \blk\blk{\idt av();\ phtg();\ av()}}\par
\vspace{0.3\baselineskip}
\vspace{\neglskip}\mbox{\blc {\mcl for}{\idt \ i\ }{\mcl in}{\idt \ range(10):}}\par
\vspace{\neglskip}\mbox{\blc \blk\blk{\idt tilde()}}\par
\vspace{\neglskip}\mbox{\blc {\idt st()}}\par

\medskip }


Si l'on pense à faire reculer le \robot\ mais que l'on oublie de le remettre
dans la bonne direction, le résultat est encore plus étrange~:

% figure et programme insecables (sinon, a un saut de page, l'espace
% negatif fait remonter le programme au-dessus de l'en-tete)
\noindent\begin{minipage}{\linewidth}
\vspace*{25mm}
\includePostScript{FriseDeTildesErr2}{0.95}{0.95}%
\vspace*{-25mm}
{
\parindent \parindentprog
\medskip

\carrenoir
\vspace{\neglskip}\mbox{\blc {\mcl from}{\idt \ LibRobot\ }{\mcl import}{\idt \ *}}\par
\vspace{0.3\baselineskip}
\vspace{\neglskip}\mbox{\blc {\cmt \#\ Dessiner\ une\ frise\ de\ tildes\ (deuxième\ erreur)}}\par
\vspace{0.3\baselineskip}
\vspace{\neglskip}\mbox{\blc {\cmt \#\ Dessiner\ un\ motif\ en\ forme\ de\ tilde\ \char126{}}}\par
\vspace{\neglskip}\mbox{\blc {\mcl def}{\idt \ tilde():}}\par
\vspace{\neglskip}\mbox{\blc \blk\blk{\idt phtd();\ av();\ phtd();\ av()}}\par
\vspace{\neglskip}\mbox{\blc \blk\blk{\idt pqtd();\ av();\ av();\ pqtg()}}\par
\vspace{\neglskip}\mbox{\blc \blk\blk{\idt av();\ phtg();\ av()}}\par
\vspace{0.3\baselineskip}
\vspace{\neglskip}\mbox{\blc {\mcl for}{\idt \ i\ }{\mcl in}{\idt \ range(13):}}\par
\vspace{\neglskip}\mbox{\blc \blk\blk{\idt tilde()}}\par
\vspace{\neglskip}\mbox{\blc \blk\blk{\cmt \#\ Reculer\ au\ début\ du\ tilde\ suivant}}\par
\vspace{\neglskip}\mbox{\blc \blk\blk{\idt lc();\ p3htg();\ av();\ bc()}}\par
\vspace{\neglskip}\mbox{\blc {\idt st()}}\par

\medskip }

\end{minipage}

Une bonne façon de trouver les erreurs dans une boucle ``for'' est donc
de chercher l'invariant, c'est-à-dire la position et l'orientation dans
lesquelles doit se trouver le \robot\ avant de dessiner le motif en
exécutant le corps de la boucle. Il faut alors vérifier que le \robot\ est
convenablement placé quand on arrive pour la première fois dans la boucle
et qu'il l'est à nouveau après avoir exécuté le corps de la boucle.

La même notion d'invariant s'applique pour les procédures. Avant chaque
appel d'une procédure, il faut que le robot soit dans un ou plusieurs
états possibles qui doivent être parfaitement définis par la position
levée ou baissée du crayon du \robot\ et par la position et l'orientation
du \robot\ par rapport à la figure dessinée par la procédure.

\noindent\rule{137mm}{0.1mm} %---------------------------------------------
\begin{exercices}{\ \\

\vspace{-\baselineskip}
\tiret Écrire des programmes \python\ utilisant des boucles ``for'' pour dessiner
quelques-unes des frises de la figure
\ref{fig:Frise-1}  page \pageref{fig:Frise-1}.

\tiret Écrire des programmes \python\ utilisant des boucles ``for'' dans des pro\-cé\-du\-res
pour dessiner quelques-unes des frises de la figure
\ref{fig:Frise-2}  page \pageref{fig:Frise-2}.

\tiret Écrire des programmes \python\ pour dessiner quelques-unes des rosaces,
frises et pavage de la figure
\ref{fig:Ro-Fr-Pa}  page \pageref{fig:Ro-Fr-Pa}, en utilisant des boucles
``for'' imbriquées indépendantes.

\tiret Écrire des programmes \python\ pour dessiner quelques-unes des spirales,
quadrillages et autres tuyaux d'orgues de la figure
\ref{fig:SpirQuadTuy}  ci-dessous, en utilisant des boucles
``for'' imbriquées dépendantes.

\tiret Écrire un programme \python\ pour dessiner le pavage de
la figure \ref{fig:PavCLC} ci-dessous en indiquant avec précision les
invariants des procédures et boucles ``for''.
\noindent $\Box$}\end{exercices}
\begin{figure}[htbp]
 \includevspacePostScript[4.5cm]{SpirQuadTuy}{0.95}{0.95}
 \caption{Quadrillages, carrés, spirale et tuyaux d'orgues}
 \label{fig:SpirQuadTuy}
\end{figure}
\begin{figure}[htbp]
 \includevspacePostScript[7.1cm]{PavageCarresTrianglesCercles}{0.95}{0.95}
 \caption{Pavage de carrés, triangles et cercles}\label{fig:PavCLC}
\end{figure}
\begin{figure}[htbp]
 \includevspacePostScript[17.8cm]{Frise-1}{0.95}{0.95}
 \caption{Frises (1)}\label{fig:Frise-1}
\end{figure}
\begin{figure}[htbp]
 \includevspacePostScript[17.8cm]{Frise-2}{0.95}{0.95}
 \caption{Frises (2)}\label{fig:Frise-2}
\end{figure}
\begin{figure}[htbp]
 \includevspacePostScript[17.8cm]{Ro-Fr-Pa}{0.95}{0.95}
 \caption{Rosaces, frises et pavage}\label{fig:Ro-Fr-Pa}
\end{figure}

\clearpage

%%%%%%%%%%%%%%%%%%%%%%%%%%%%%%%%%%%%%%%%%%%%%%%%%%%%%%%%%%%%%%%%%%%%%%%%%%%%%%%%%%%%%%%%%%%%%%
\begin{corriges}{
\renewcommand{\carrenoir}{\hspace{-4mm}\rule{1.5mm}{1.5mm}\hspace{2.5mm}}
Nous donnons quelques programmes, les autres se trouvant dans le dossier
{\idt livre/\allowbreak 07\_for}~:
{
\parindent \parindentprog
\medskip

\carrenoir
\vspace{\neglskip}\mbox{\blc {\mcl from}{\idt \ LibRobot\ }{\mcl import}{\idt \ *}}\par
\vspace{0.3\baselineskip}
\vspace{\neglskip}\mbox{\blc {\cmt \#\ Dessiner\ une\ frise\ crantée}}\par
\vspace{0.3\baselineskip}
\vspace{\neglskip}\mbox{\blc {\cmt \#\ Dessiner\ le\ motif\ \textquotesingle{}n\_\textquotesingle{}}}\par
\vspace{\neglskip}\mbox{\blc {\mcl def}{\idt \ cran():}}\par
\vspace{\neglskip}\mbox{\blc \blk\blk{\idt av();\ pqtd();\ av();\ pqtd();\ av();\ pqtg();\ av()}}\par
\vspace{\neglskip}\mbox{\blc \blk\blk{\idt pqtg()}}\par
\vspace{0.3\baselineskip}
\vspace{\neglskip}\mbox{\blc {\idt ag()}}\par
\vspace{\neglskip}\mbox{\blc {\mcl for}{\idt \ i\ }{\mcl in}{\idt \ range(20):}}\par
\vspace{\neglskip}\mbox{\blc \blk\blk{\idt cran()}}\par
\vspace{\neglskip}\mbox{\blc {\idt st()}}\par

\medskip }

{
\parindent \parindentprog
\medskip

\carrenoir
\vspace{\neglskip}\mbox{\blc {\mcl from}{\idt \ LibRobot\ }{\mcl import}{\idt \ *}}\par
\vspace{0.3\baselineskip}
\vspace{\neglskip}\mbox{\blc {\cmt \#\ Dessiner\ un\ quadrillage\ triangulaire}}\par
\vspace{\neglskip}\mbox{\blc {\cmt \#\ Dessiner\ 10\ lignes\ de\ carrés\ juxtaposés}}\par
\vspace{\neglskip}\mbox{\blc {\cmt \#\ horizontalement}}\par
\vspace{\neglskip}\mbox{\blc {\idt pqtd()}}\par
\vspace{\neglskip}\mbox{\blc {\mcl for}{\idt \ i\ }{\mcl in}{\idt \ range(1,\ 11):}}\par
\vspace{\neglskip}\mbox{\blc \blk\blk{\cmt \#\ Dessiner\ une\ ligne\ de\ i\ carrés\ juxtaposés}}\par
\vspace{\neglskip}\mbox{\blc \blk\blk{\cmt \#\ horizontalement,\ en\ partant\ en\ haut\ à\ gauche\ et}}\par
\vspace{\neglskip}\mbox{\blc \blk\blk{\cmt \#\ en\ arrivant\ en\ haut\ à\ droite,\ le\ robot\ étant}}\par
\vspace{\neglskip}\mbox{\blc \blk\blk{\cmt \#\ dans\ les\ deux\ cas\ orienté\ vers\ l\textquotesingle{}est.}}\par
\vspace{\neglskip}\mbox{\blc \blk\blk{\mcl for}{\idt \ j\ }{\mcl in}{\idt \ range(i):}}\par
\vspace{\neglskip}\mbox{\blc \blk\blk\blk\blk{\cmt \#\ Dessiner\ un\ carré,\ en\ partant\ du\ coin}}\par
\vspace{\neglskip}\mbox{\blc \blk\blk\blk\blk{\cmt \#\ supérieur\ gauche\ et\ en\ arrivant\ au\ coin}}\par
\vspace{\neglskip}\mbox{\blc \blk\blk\blk\blk{\cmt \#\ supérieur\ droit,\ le\ robot\ étant\ dans\ les\ deux}}\par
\vspace{\neglskip}\mbox{\blc \blk\blk\blk\blk{\cmt \#\ cas\ orienté\ vers\ l\textquotesingle{}est.}}\par
\vspace{\neglskip}\mbox{\blc \blk\blk\blk\blk{\mcl for}{\idt \ k\ }{\mcl in}{\idt \ range(4):}}\par
\vspace{\neglskip}\mbox{\blc \blk\blk\blk\blk\blk\blk{\idt av();\ pqtd()}}\par
\vspace{\neglskip}\mbox{\blc \blk\blk\blk\blk{\idt av()}}\par
\vspace{\neglskip}\mbox{\blc \blk\blk{\cmt \#\ Revenir,\ sans\ dessiner,\ au\ début\ de\ la\ ligne}}\par
\vspace{\neglskip}\mbox{\blc \blk\blk{\cmt \#\ suivante,\ le\ robot\ étant\ orienté\ vers\ l\textquotesingle{}est\ au}}\par
\vspace{\neglskip}\mbox{\blc \blk\blk{\cmt \#\ départ\ et\ à\ l\textquotesingle{}arrivée.}}\par
\vspace{\neglskip}\mbox{\blc \blk\blk{\idt lc();\ pdt()}}\par
\vspace{\neglskip}\mbox{\blc \blk\blk{\mcl for}{\idt \ j\ }{\mcl in}{\idt \ range(i):}}\par
\vspace{\neglskip}\mbox{\blc \blk\blk\blk\blk{\idt av()}}\par
\vspace{\neglskip}\mbox{\blc \blk\blk{\idt pqtg();\ av();\ pqtg();\ bc()}}\par
\vspace{\neglskip}\mbox{\blc {\idt st()}}\par

\medskip }

{
\parindent \parindentprog
\medskip

\carrenoir
\vspace{\neglskip}\mbox{\blc {\mcl from}{\idt \ LibRobot\ }{\mcl import}{\idt \ *}}\par
\vspace{0.3\baselineskip}
\vspace{\neglskip}\mbox{\blc {\cmt \#\ Dessiner\ un\ pavage\ de\ carrés,\ triangles\ et\ cercles}}\par
\vspace{0.3\baselineskip}
\vspace{\neglskip}\mbox{\blc {\cmt \#\ Dessiner\ un\ carré\ de\ côté\ de\ taille\ 2,\ le\ robot}}\par
\vspace{\neglskip}\mbox{\blc {\cmt \#\ partant\ du\ coin\ inférieur\ gauche\ tourné\ vers\ l\textquotesingle{}est,}}\par
\vspace{\neglskip}\mbox{\blc {\cmt \#\ et\ arrivant\ crayon\ baissé\ au\ coin\ inférieur\ droit}}\par
\vspace{\neglskip}\mbox{\blc {\cmt \#\ également\ orienté\ à\ l\textquotesingle{}est.}}\par
\vspace{\neglskip}\mbox{\blc {\mcl def}{\idt \ carre():}}\par
\vspace{\neglskip}\mbox{\blc \blk\blk{\cmt \#\ Dessiner\ le\ carré.}}\par
\vspace{\neglskip}\mbox{\blc \blk\blk{\idt av();\ av();\ pqtg();\ av();\ av();\ pqtg()}}\par
\vspace{\neglskip}\mbox{\blc \blk\blk{\idt av();\ av();\ pqtg();\ av();\ av();\ pqtg()}}\par
\vspace{\neglskip}\mbox{\blc \blk\blk{\cmt \#\ Avancer\ au\ coin\ inférieur\ droit.}}\par
\vspace{\neglskip}\mbox{\blc \blk\blk{\idt av();\ av()}}\par
\vspace{0.3\baselineskip}
\vspace{\neglskip}\mbox{\blc {\cmt \#\ Dessiner\ deux\ triangles\ opposés\ par\ un\ sommet}}\par
\vspace{\neglskip}\mbox{\blc {\cmt \#\ inscrits\ dans\ un\ carré\ invisible\ du\ quadrillage\ de}}\par
\vspace{\neglskip}\mbox{\blc {\cmt \#\ côté\ 2,\ en\ partant\ du\ coin\ inférieur\ gauche\ du\ carré}}\par
\vspace{\neglskip}\mbox{\blc {\cmt \#\ tourné\ vers\ l\textquotesingle{}est,\ et\ en\ arrivant\ au\ coin\ inférieur}}\par
\vspace{\neglskip}\mbox{\blc {\cmt \#\ droit,\ crayon\ baissé\ et\ tourné\ vers\ l\textquotesingle{}est.}}\par
\vspace{\neglskip}\mbox{\blc {\mcl def}{\idt \ triangles():}}\par
\vspace{\neglskip}\mbox{\blc \blk\blk{\idt lc();\ av();\ bc()}}\par
\vspace{\neglskip}\mbox{\blc \blk\blk{\idt pqtg();\ av();\ av();\ p3htg();\ av();\ p3htg()}}\par
\vspace{\neglskip}\mbox{\blc \blk\blk{\idt av();\ av();\ p3htd();\ av();\ p3htg()}}\par
\vspace{\neglskip}\mbox{\blc \blk\blk{\idt lc();\ av();\ bc()}}\par
\vspace{0.3\baselineskip}
\vspace{\neglskip}\mbox{\blc {\cmt \#\ Dessiner\ un\ cercle\ de\ rayon\ 1\ inscrit\ dans\ un\ carré}}\par
\vspace{\neglskip}\mbox{\blc {\cmt \#\ invisible\ du\ quadrillage\ de\ côté\ 2,\ en\ partant\ du}}\par
\vspace{\neglskip}\mbox{\blc {\cmt \#\ coin\ inférieur\ gauche\ du\ carré\ tourné\ vers\ l\textquotesingle{}est,\ et}}\par
\vspace{\neglskip}\mbox{\blc {\cmt \#\ en\ arrivant\ au\ coin\ inférieur\ droit\ du\ carré,\ crayon}}\par
\vspace{\neglskip}\mbox{\blc {\cmt \#\ baissé\ et\ tourné\ vers\ l\textquotesingle{}est.}}\par
\vspace{\neglskip}\mbox{\blc {\mcl def}{\idt \ cercle():}}\par
\vspace{\neglskip}\mbox{\blc \blk\blk{\idt lc();\ av();\ bc()}}\par
\vspace{\neglskip}\mbox{\blc \blk\blk{\idt vg();\ vg();\ vg();\ vg()}}\par
\vspace{\neglskip}\mbox{\blc \blk\blk{\idt lc();\ av();\ bc()}}\par
\vspace{0.3\baselineskip}
\vspace{\neglskip}\mbox{\blc {\cmt \#\ Dessiner\ une\ ligne\ constituée\ d\textquotesingle{}un\ motif\ répété\ 7}}\par
\vspace{\neglskip}\mbox{\blc {\cmt \#\ fois,\ le\ motif\ étant\ formé\ d\textquotesingle{}un\ carré,\ de\ triangles}}\par
\vspace{\neglskip}\mbox{\blc {\cmt \#\ et\ d\textquotesingle{}un\ cercle\ juxtaposés,\ le\ robot\ partant\ du\ coin}}\par
\vspace{\neglskip}\mbox{\blc {\cmt \#\ inférieur\ gauche\ de\ la\ ligne\ orienté\ vers\ l\textquotesingle{}est,\ et}}\par
\vspace{\neglskip}\mbox{\blc {\cmt \#\ arrivant\ au\ coin\ inférieur\ droit\ également\ orienté\ à}}\par
\vspace{\neglskip}\mbox{\blc {\cmt \#\ l\textquotesingle{}est.}}\par
\vspace{\neglskip}\mbox{\blc {\mcl def}{\idt \ ligneCarre():}}\par
\vspace{\neglskip}\mbox{\blc \blk\blk{\mcl for}{\idt \ i\ }{\mcl in}{\idt \ range(7):}}\par
\vspace{\neglskip}\mbox{\blc \blk\blk\blk\blk{\cmt \#\ Le\ robot\ est\ crayon\ baissé,\ à\ l\textquotesingle{}est,\ en\ bas\ à}}\par
\vspace{\neglskip}\mbox{\blc \blk\blk\blk\blk{\cmt \#\ gauche\ du\ motif}}\par
\vspace{\neglskip}\mbox{\blc \blk\blk\blk\blk{\idt carre();\ triangles();\ cercle()}}\par
\vspace{0.3\baselineskip}
\vspace{\neglskip}\mbox{\blc {\cmt \#\ Comme\ ligneCarre,\ avec\ le\ motif\ formé\ de\ triangles,}}\par
\vspace{\neglskip}\mbox{\blc {\cmt \#\ d\textquotesingle{}un\ cercle\ et\ d\textquotesingle{}un\ carré.}}\par
\vspace{\neglskip}\mbox{\blc {\mcl def}{\idt \ ligneTriangles():}}\par
\vspace{\neglskip}\mbox{\blc \blk\blk{\mcl for}{\idt \ i\ }{\mcl in}{\idt \ range(7):}}\par
\vspace{\neglskip}\mbox{\blc \blk\blk\blk\blk{\idt triangles();\ cercle();\ carre()}}\par
\vspace{0.3\baselineskip}
\vspace{\neglskip}\mbox{\blc {\cmt \#\ Comme\ ligneCarre,\ avec\ le\ motif\ formé\ d\textquotesingle{}un\ cercle,}}\par
\vspace{\neglskip}\mbox{\blc {\cmt \#\ d\textquotesingle{}un\ carré\ et\ de\ triangles.}}\par
\vspace{\neglskip}\mbox{\blc {\mcl def}{\idt \ ligneCercle():}}\par
\vspace{\neglskip}\mbox{\blc \blk\blk{\mcl for}{\idt \ i\ }{\mcl in}{\idt \ range(7):}}\par
\vspace{\neglskip}\mbox{\blc \blk\blk\blk\blk{\idt cercle();\ carre();\ triangles()}}\par
\vspace{0.3\baselineskip}
\vspace{\neglskip}\mbox{\blc {\cmt \#\ Ramener\ le\ robot\ du\ coin\ inférieur\ droit\ d\textquotesingle{}une\ ligne}}\par
\vspace{\neglskip}\mbox{\blc {\cmt \#\ au\ coin\ inférieur\ gauche\ de\ la\ ligne\ suivante,\ sans}}\par
\vspace{\neglskip}\mbox{\blc {\cmt \#\ dessiner,\ orienté\ à\ l\textquotesingle{}est\ et\ crayon\ baissé.}}\par
\vspace{\neglskip}\mbox{\blc {\mcl def}{\idt \ alaLigne():}}\par
\vspace{\neglskip}\mbox{\blc \blk\blk{\idt lc();\ ag();\ pqtd();\ av();\ av();\ pqtg();\ bc()}}\par
\vspace{0.3\baselineskip}
\vspace{\neglskip}\mbox{\blc {\cmt \#\ Paver\ l\textquotesingle{}écran\ de\ 11\ lignes\ de\ motifs\ commençant\ par}}\par
\vspace{\neglskip}\mbox{\blc {\cmt \#\ un\ carré,\ puis\ par\ des\ triangles,\ puis\ par\ un}}\par
\vspace{\neglskip}\mbox{\blc {\cmt \#\ cercle,\ du\ coin\ supérieur\ droit\ au\ coin\ inférieur}}\par
\vspace{\neglskip}\mbox{\blc {\cmt \#\ droit\ de\ l\textquotesingle{}écran,\ le\ robot\ étant\ orienté\ à\ l\textquotesingle{}est.}}\par
\vspace{\neglskip}\mbox{\blc {\mcl def}{\idt \ pavage():}}\par
\vspace{\neglskip}\mbox{\blc \blk\blk{\cmt \#\ Dessiner\ les\ lignes\ 1\ à\ 9\ de\ haut\ en\ bas.}}\par
\vspace{\neglskip}\mbox{\blc \blk\blk{\mcl for}{\idt \ j\ }{\mcl in}{\idt \ range(3):}}\par
\vspace{\neglskip}\mbox{\blc \blk\blk\blk\blk{\idt alaLigne();\ ligneCarre()}}\par
\vspace{\neglskip}\mbox{\blc \blk\blk\blk\blk{\idt alaLigne();\ ligneTriangles()}}\par
\vspace{\neglskip}\mbox{\blc \blk\blk\blk\blk{\idt alaLigne();\ ligneCercle()}}\par
\vspace{\neglskip}\mbox{\blc \blk\blk{\cmt \#\ Dessiner\ les\ lignes\ 10\ puis\ 11.}}\par
\vspace{\neglskip}\mbox{\blc \blk\blk{\idt alaLigne();\ ligneCarre()}}\par
\vspace{\neglskip}\mbox{\blc \blk\blk{\idt alaLigne();\ ligneTriangles()}}\par
\vspace{0.3\baselineskip}
\vspace{\neglskip}\mbox{\blc {\cmt \#\ Aller\ au\ coin\ supérieur\ droit,\ orienté\ à\ l\textquotesingle{}est.}}\par
\vspace{\neglskip}\mbox{\blc {\idt eh();\ ad();\ pqtd()}}\par
\vspace{\neglskip}\mbox{\blc {\cmt \#\ Paver\ l\textquotesingle{}écran.}}\par
\vspace{\neglskip}\mbox{\blc {\idt pavage();\ st()}}\par

\medskip }

\noindent $\Box$\vfill}\end{corriges}

%%%%%%%%%%%%%%%%%%%%%%%%%%%%%%%%%%%%%%%%%%%%%%%%%%%%%%%%%%%%%%%%%%%%%%%%%%%
\noindent \framebox[\textwidth]{
\begin{minipage}{\myparboxwidth}
\setlength{\parindent}{2em}
\small \sf
\vspace{\abovedisplayskip}
\section*{Addenda}

La boucle ``for'' est une des premières notions linguistiques qui soit
apparue en informatique puisqu'elle fût inventée à l'ETH de Zürich
 en 1952 par Heinz Rutishauser \cite{rutishauser}.

La comparaison des dessins que l'on peut obtenir en utilisant la boucle
``for'' avec ceux que l'on peut faire en mode interactif donne la mesure
de la puissance des ordinateurs quand il s'agit d'effectuer rapidement des
tâches répétitives. En cela l'ordinateur se distingue très
nettement des calculettes non programmables. On comprend également
que cette puissance de calcul n'a rien à voir avec une quelconque forme
de pensée ou d'intelligence si ce n'est celle du programmeur de
l'ordinateur.

En \python, on utilise une boucle {\mcl for} {\idt i} {\mcl in} {\idt range(}$n_1${\idt ,\ }$n_2${\idt ):}
 \ldots\ pour répéter
$n_2 - n_1$ fois le corps de la boucle (les instructions décalées qui suivent
les deux-points) sauf si $n_2 \leq n_1$ auquel cas le corps de la boucle n'est pas
exécuté. A chaque itération, le compteur de boucle {\idt i} prend les
valeurs successives $n_1$, $n_1 + 1$, $n_1 + 2$, \ldots, $n_2 - 1$. Dans la forme {\mcl for}
{\idt i} {\mcl in} {\idt range(}$n_1${\idt ,\ }$n_2${\idt ,\ -1):} \ldots,
{\idt i} prend les valeurs successives $n_1$, $n_1 - 1$, $n_1 - 2$, \ldots, $n_2 + 1$
sauf si $n_2 \geq n_1$ auquel cas la boucle n'est pas exécutée~; par exemple
{\mcl for} {\idt i} {\mcl in} {\idt range(}$n${\idt ,\ 0,\ -1):} fait prendre à {\idt i} les
valeurs $n$, $n-1$, \ldots, 1 (le troisième nombre, ici $-1$, est le pas
qui est ajouté au compteur à chaque itération). Si le corps de la boucle modifie la
valeur du compteur, cela n'a pas d'effet sur les itérations suivantes~: à chaque
itération, le compteur reprend la valeur suivante de la suite {\idt range}.
Quand le compteur ne sert à rien dans le corps de la boucle, les programmeurs
\python\ le nomment souvent `\_', comme dans
{\mcl for}\ {\idt \_\ }{\mcl in}\ {\idt range(16):}\ \ldots

L'utilisation de boucles ``for'' exerce la capacité de percevoir
quand un dessin peut se décomposer en un certain nombre de sous-dessins
identiques (dans le cas d'une boucle simple) ou similaires (dans le cas
de boucles imbriquées). C'est un exercice quotidien pour l'informaticien
qui doit savoir décomposer une tâche complexe en sous-tâches simples,
ce qui demande parfois énormément d'imagination.
Il s'agit d'un premier pas vers l'acquisition du raisonnement
inductif des scien\-ti\-fi\-ques qui, à partir d'observations spécifiques,
créent une hypothèse générale rendant compte de ces observations,
hypothèse qu'il reste à valider dans tous les cas possibles, ici de
manière tout à fait expérimentale, en observant le dessin à
l'écran. Les mathématiciens appellent {\em conjecture} une telle
hypothèse induite de l'expérience et non encore démontrée.

\vspace{\abovedisplayskip}
\end{minipage}
}
%%%%%%%%%%%%%%%%%%%%%%%%%%%%%%%%%%%%%%%%%%%%%%%%%%%%%%%%%%%%%%%%%%%%%%%%%%%%
%%%%%%%%%%%%%%%%%%%%%%%%%%%%%%%%%%%%%%%%%%%%%%%%%%%%%%%%%%%%%%%%%%%%%%%%%%%
\noindent \framebox[\textwidth]{
\begin{minipage}{\myparboxwidth}
\setlength{\parindent}{2em}
\small \sf
\vspace{\abovedisplayskip}
Ce raisonnement inductif a longtemps été confondu avec le raisonnement
par récurrence qui est essentiel en mathématiques. Il a fallu plus d'un
millénaire pour que les mathématiciens
(Fermat\footnote{\ \href{https://en.wikipedia.org/wiki/Pierre_de_Fermat}{Pierre de Fermat}, 1601--1665, mathématicien français.},
Pascal\footnote{\ \href{https://en.wikipedia.org/wiki/Blaise_Pascal}{Blaise Pascal}, 1623--1662, mathématicien français.}
notamment) comprennent
pleinement ce raisonnement par récurrence (\cite{knuth2} page 17).

Nous en
sommes pour l'instant aux prémisses, l'absence d'induction correcte
 étant encore admise dans le raisonnement
par récurrence~: si le dessin est bien
fait pour le motif de base et pour
la première répétition du motif alors il est réputé correct, seul le
problème du cadrage dans la fenêtre restant à résoudre par simple
expérimentation.

La notion d'{\em invariant} (c'est-à-dire pour notre \robot\ la
description de ses positions et orientations possibles au début du
corps des boucles) est due à Turing
\footnote{\ \href{https://en.wikipedia.org/wiki/Alan_Turing}{Alan Turing}, informaticien anglais.} en 1949, oubliée, puis redécouverte pat
Naur%
\footnote{\ \href{https://en.wikipedia.org/wiki/Peter_Naur}{Peter Naur}, informaticien danois.} en 1966 et Floyd%
\footnote{\ \href{https://en.wikipedia.org/wiki/Robert_W._Floyd}{Robert Floyd}, informaticien américain.} en 1967.
Elle repose
sur un raisonnement par récurrence~: si l'invariant est vrai quand on
arrive pour la première fois dans la boucle (c'est-à-dire que la
position et l'orientation du \robot\ est conforme à la description qui en
est faite dans l'invariant) et si, en supposant que l'invariant soit vrai au
début de l'exécution du corps de la boucle, il reste encore vrai après
l'exécution de ce corps de boucle alors l'invariant est toujours vrai au
début de chaque itération dans la boucle et l'est encore quand
on sort de la boucle. La notion s'applique aussi bien aux procédures~: on
peut utiliser un
{\em invariant d'entrée} décrivant les positions et orientations du
\robot\ avant l'appel et un {\em invariant de sortie} décrivant
ces positions et orientations après l'appel, de sorte que le corps de la
procédure soit compréhensible sans avoir à consulter tous les points
d'appel pour savoir dans quelles conditions la procédure est appelée.

\vspace{\abovedisplayskip}
\end{minipage}
}
%%%%%%%%%%%%%%%%%%%%%%%%%%%%%%%%%%%%%%%%%%%%%%%%%%%%%%%%%%%%%%%%%%%%%%%%%%%%
%%%%%%%%%%%%%%%%%%%%%%%%%%%%%%%%%%%%%%%%%%%%%%%%%%%%%%%%%%%%%%%%%%%%%%%%%%%
\noindent \framebox[\textwidth]{
\begin{minipage}{\myparboxwidth}
\setlength{\parindent}{2em}
\small \sf
\vspace{\abovedisplayskip}
L'ordinateur peut être utilisé pour manipuler des programmes, en
particulier pour découvrir automatiquement des invariants de forme
simple en chaque point d'un programme. C'est l'objet de
l'{\em interprétation abstraite} des programmes
\cite{CousotIA}. Pour le \robot, par l'exemple, une forme simple
d'invariant est constituée par l'ensemble de ses orientations possibles,
en ignorant la position de son crayon, ses coordonnées sur
l'écran, etc. Les orientations possibles du \robot\ concret peuvent être
calculées en exécutant le programme avec un \robot\ abstrait qui
exécute les commandes du \robot\ concret sans bouger
mais en changeant simplement d'orientation. Pour ce \robot\ abstrait
la commande {\idt av} ne fait rien car elle ne change pas l'orientation du
\robot\ concret$\:$; la commande {\idt vd} a le même effet que
{\idt pqtd}$\:$; etc. Pendant l'exécution abstraite les orientations
possibles du \robot\ abstrait sont notées comme commentaire
en chaque point du programme. L'exécution abstraite du programme
s'arrête après stabilisation des résultats. Par exemple, pour le programme {\idt FriseDeN\_2} (page
\pageref{prog:FriseDeN_2}), on trouve~:

{
\parindent \parindentprog
\medskip

\carrenoir
\vspace{\neglskip}\mbox{\blc {\mcl from}{\idt \ LibRobot\ }{\mcl import}{\idt \ *}}\par
\vspace{0.3\baselineskip}
\vspace{\neglskip}\mbox{\blc {\cmt \#\ Dessiner\ une\ frise\ de\ N\ avec\ une\ boucle\ \char34{}for\char34{}}}\par
\vspace{\neglskip}\mbox{\blc {\cmt \#\ Nord}}\par
\vspace{\neglskip}\mbox{\blc {\idt ag()}}\par
\vspace{\neglskip}\mbox{\blc {\cmt \#\ Nord}}\par
\vspace{\neglskip}\mbox{\blc {\mcl for}{\idt \ j\ }{\mcl in}{\idt \ range(16):}}\par
\vspace{\neglskip}\mbox{\blc \blk\blk{\idt av();\ av()\ \ }{\cmt \#\ Nord\ \ \#\ Nord}}\par
\vspace{\neglskip}\mbox{\blc \blk\blk{\idt pdt();\ vg()\ \ }{\cmt \#\ Nord\ \ \#\ Sud}}\par
\vspace{\neglskip}\mbox{\blc \blk\blk{\idt vd();\ pdt()\ \ }{\cmt \#\ Est\ \ \#\ Sud}}\par
\vspace{\neglskip}\mbox{\blc {\cmt \#\ Nord}}\par
\vspace{\neglskip}\mbox{\blc {\cmt \#\ Nord}}\par
\vspace{\neglskip}\mbox{\blc {\idt st()}}\par
\vspace{\neglskip}\mbox{\blc {\cmt \#\ Nord}}\par

\medskip }


La règle des signes est un exemple d'interprétation abstraite utilisée
en mathématiques. Le signe d'une expression peut être
calculé en utilisant le signe des variables, qui remplace leurs valeurs,
grâce aux opérations abstraites comme $\oplus + \oplus = \oplus$,
$\oplus - \ominus = \oplus$, $\oplus \times \ominus = \ominus$, etc. Le prix
à payer pour cette simplification est évidemment l'imprécision des
résultats, comme pour $\oplus + \ominus$.
\vspace{\abovedisplayskip}
\end{minipage}
}
%%%%%%%%%%%%%%%%%%%%%%%%%%%%%%%%%%%%%%%%%%%%%%%%%%%%%%%%%%%%%%%%%%%%%%%%%%%%
